\documentclass[a4paper,12pt]{article}
\usepackage[T2A]{fontenc}
\usepackage[utf8]{inputenc}
\usepackage[russian]{babel}
\usepackage{amsmath,amssymb,mathtools}
\usepackage{PTSans}
\renewcommand{\familydefault}{\sfdefault}
\usepackage{icomma}
\usepackage{geometry}
\geometry{left=18mm,right=18mm,top=18mm,bottom=20mm}
\usepackage{microtype}
\usepackage{tikz}
\usetikzlibrary{calc,arrows.meta,positioning,shapes.geometric}
\usepackage{tkz-euclide}
\usepackage{tabularx,booktabs,longtable}
\usepackage{enumitem}
\usepackage{hyperref}
\usepackage{parskip}
\usepackage{fancyhdr}
\usepackage{setspace}
\usepackage{xcolor}

\definecolor{accent}{HTML}{356B75}
\definecolor{darkaccent}{HTML}{234B52}
\definecolor{textgray}{HTML}{333333}
\definecolor{soft}{HTML}{EEF4F5}
\definecolor{muted}{HTML}{6F7C80}

\hypersetup{colorlinks=true,linkcolor=darkaccent,urlcolor=darkaccent}
\setlength{\parindent}{0pt}
\setlength{\headheight}{25pt}
\onehalfspacing
\setlist[itemize]{leftmargin=6mm,itemsep=2pt,topsep=3pt}
\setlist[enumerate]{leftmargin=8mm,itemsep=4pt,topsep=3pt}
\pagestyle{fancy}
\fancyhf{}
\fancyhead[L]{\small\color{muted}Чек-лист: оптимизационные задачи}
\fancyhead[R]{\small\color{muted}28.09.26}
\fancyfoot[C]{\small\thepage}
\renewcommand{\headrulewidth}{0.3pt}
\renewcommand{\headrule}{\hbox to\headwidth{\color{accent}\leaders\hrule height \headrulewidth\hfill}}

\newcommand{\sectionline}{\vspace{1mm}{\color{accent}\hrule height 0.7pt}\vspace{3mm}}
\newcommand{\key}[1]{\textbf{\color{darkaccent}#1}}

\begin{document}

\thispagestyle{empty}
\begin{center}
\vspace*{20mm}
{\Large\color{muted}Математика · профильный уровень}\par
\vspace{11mm}
{\fontsize{30}{34}\selectfont\bfseries\color{darkaccent}Чек-лист}\par
\vspace{4mm}
{\LARGE\bfseries Оптимизационные задачи с производной}\par
\vspace{10mm}
{\large Как переводить прикладной текст в функцию одной переменной, учитывать ограничения и доказывать максимум или минимум}\par
\vspace{18mm}
{\Large\bfseries Николь Саркисьянц}\par
\vspace{5mm}
{\large 28.09.26}\par
\vfill
{\small\color{textgray}Лёвин Артём Александрович эксклюзивно для Николь Саркисьянц}\par
\vspace{8mm}
\begin{tikzpicture}[x=1cm,y=1cm]
\draw[accent,line width=1.2pt] (-7,0) .. controls (-3,1.0) and (2,-1.0) .. (7,0.2);
\draw[accent!50,line width=0.7pt] (-6.4,-0.25) .. controls (-2.5,0.55) and (2.8,-0.7) .. (6.3,-0.1);
\end{tikzpicture}
\end{center}
\clearpage

\section*{1. Главная идея}
\sectionline
В прикладной задаче слова \key{«наибольший»}, \key{«наименьший»}, \key{«максимальный»}, \key{«минимальный»} указывают, \textbf{какую величину нужно оптимизировать}. Для этой величины строят функцию одной переменной, а затем исследуют её на допустимых значениях с помощью производной.

\key{Главный вопрос при чтении условия:} \emph{к чему относится слово «максимальный» или «минимальный»?}

\begin{itemize}
  \item максимальная выручка $\Rightarrow$ строим функцию выручки $R$;
  \item максимальная площадь $\Rightarrow$ строим функцию площади $S$;
  \item максимальная мощность $\Rightarrow$ строим функцию мощности $P$;
  \item максимальный объём $\Rightarrow$ строим функцию объёма $V$.
\end{itemize}

\section*{2. Универсальный маршрут решения}
\sectionline
\begin{enumerate}
  \item \key{Обозначьте переменную}, которую требуется найти, и подпишите её смысл.
  \item \key{Найдите оптимизируемую величину}: что должно стать максимальным или минимальным?
  \item \key{Составьте формулу} этой величины из данных условия.
  \item \key{Сведите формулу к одной переменной}. Желательно к той, которую спрашивают в задаче.
  \item \key{Запишите допустимые значения}: длины, количества, температуры, сопротивления и т.~п.
  \item \key{Упростите функцию}, если это делает дифференцирование прозрачнее.
  \item \key{Найдите производную} и решите $f'(x)=0$ внутри допустимой области.
  \item \key{Исследуйте знак производной} только на допустимом промежутке.
  \item Если границы входят в область определения, \key{не забудьте их}: при необходимости сравните значения функции в критических и граничных точках.
  \item \key{Ответьте ровно на вопрос}: найдите требуемую переменную и, если просят, само максимальное или минимальное значение.
\end{enumerate}

\section*{3. Ограничения --- часть решения}
\sectionline
В прикладных задачах алгебраическая формула может существовать при значениях, которые не имеют смысла в рассматриваемой ситуации.

\begin{itemize}
  \item количество товара: обычно $q\in\mathbb N$; для применения производной сначала исследуют непрерывную модель, а затем возвращаются к целым значениям;
  \item длина, площадь, объём и сопротивление в рассматриваемых моделях должны удовлетворять физическим ограничениям задачи;
  \item если из квадрата со стороной $30$ вырезают по углам квадраты со стороной $x$, то
  \[
    0<x<15,
  \]
  поскольку сторона дна $30-2x$ должна быть положительной;
  \item если условием дан промежуток, например $20\le t\le80$, исследование ведётся только на нём.
\end{itemize}

\textbf{Для целочисленной величины:} если непрерывная модель даёт нецелую точку максимума или минимума, сравните ближайшие допустимые целые значения.

\textbf{Важно:} критическая точка вне допустимой области не может быть ответом прикладной задачи.

\clearpage
\section*{Алгоритм: как решать оптимизационную задачу}
\vspace{4mm}
\begin{center}
\begin{tikzpicture}[
  node distance=6mm and 13mm,
  every node/.style={font=\small,align=center},
  startstop/.style={ellipse,draw=accent,line width=0.9pt,minimum width=42mm,minimum height=9mm,fill=soft},
  process/.style={rectangle,rounded corners=2mm,draw=accent,line width=0.9pt,text width=101mm,minimum height=10mm,fill=white},
  decision/.style={diamond,aspect=2.2,draw=accent,line width=0.9pt,text width=45mm,inner sep=1.5mm,fill=soft},
  side/.style={rectangle,rounded corners=2mm,draw=accent,line width=0.9pt,text width=47mm,minimum height=10mm,fill=white},
  arrow/.style={-{Stealth[length=2.6mm]},line width=0.9pt,color=darkaccent}
]
\node[startstop] (start) {Прочитать условие};
\node[process,below=of start] (target) {Определить, какая величина должна быть максимальной или минимальной};
\node[process,below=of target] (formula) {Составить целевую функцию};
\node[decision,below=of formula] (onevar) {Осталась одна переменная?};
\node[side,right=19mm of onevar] (sub) {Выразить лишние величины и подставить};
\node[process,below=of onevar] (domain) {Записать допустимую область и смысловые ограничения};
\node[process,below=of domain] (der) {Найти $f'(x)$ и упростить её до удобной для анализа формы};
\node[process,below=of der] (crit) {Найти критические точки внутри допустимой области};
\node[process,below=of crit] (signs) {Расставить знаки $f'(x)$ на допустимых промежутках};
\node[process,below=of signs] (bounds) {Учесть допустимые границы, если они входят в область};
\node[process,below=of bounds] (proof) {Обосновать максимум или минимум по знакам производной и, когда нужно, по значениям функции};
\node[startstop,below=of proof] (answer) {Записать ответ на вопрос задачи};

\draw[arrow] (start)--(target);
\draw[arrow] (target)--(formula);
\draw[arrow] (formula)--(onevar);
\draw[arrow] (onevar)-- node[left]{да} (domain);
\draw[arrow] (onevar)-- node[above]{нет} (sub);
\draw[arrow] (sub) |- (formula.east);
\draw[arrow] (domain)--(der);
\draw[arrow] (der)--(crit);
\draw[arrow] (crit)--(signs);
\draw[arrow] (signs)--(bounds);
\draw[arrow] (bounds)--(proof);
\draw[arrow] (proof)--(answer);
\end{tikzpicture}
\end{center}
\clearpage

\section*{4. Типовой пример: максимальная выручка}
\sectionline
Пусть $q$ --- количество проданных наборов, а цена одного набора задаётся формулой
\[
  p(q)=180-3q.
\]
Выручка равна цене, умноженной на количество:
\[
  R(q)=p(q)\,q=(180-3q)q=180q-3q^2.
\]
По смыслу $q\in\mathbb N$, $q>0$, а цена должна быть положительной: $180-3q>0$, то есть $q<60$.

Для поиска вершины сначала исследуем непрерывную модель на $0<q<60$:
\[
  R'(q)=180-6q.
\]
Критическая точка:
\[
  180-6q=0 \quad\Longrightarrow\quad q=30.
\]

\begin{center}
\begin{tikzpicture}[x=1.15cm,y=1cm,font=\small]
\draw[-{Stealth[length=2mm]},line width=0.8pt] (0,0)--(8,0) node[right]{$q$};
\draw[line width=0.8pt] (3.5,-0.13)--(3.5,0.13) node[below=4pt]{$30$};
\node at (1.8,0.45) {$R'(q)>0$};
\node at (5.4,0.45) {$R'(q)<0$};
\draw[-{Stealth[length=2mm]},accent,line width=1pt] (1,0.9)--(3.15,1.45);
\draw[-{Stealth[length=2mm]},accent,line width=1pt] (3.85,1.45)--(6.1,0.9);
\node[darkaccent] at (3.5,1.72) {максимум};
\end{tikzpicture}
\end{center}

До $30$ выручка возрастает, после $30$ убывает. Точка $q=30$ уже целая, поэтому дополнительного сравнения соседних целых значений не требуется. Следовательно,
\[
  q=30,\qquad R_{\max}=R(30)=(180-90)\cdot30=2700.
\]

\section*{5. Если целевая функция --- произведение}
\sectionline
Пусть
\[
  A=t-20,\qquad B=80-t,\qquad Q=AB,\qquad 20\le t\le80.
\]
Тогда
\[
  Q(t)=(t-20)(80-t).
\]
Можно раскрыть скобки или применить правило произведения:
\[
  (uv)'=u'v+uv'.
\]
Получаем
\[
  Q'(t)=(80-t)-(t-20)=100-2t.
\]
Отсюда $t=50$. При $20\le t<50$ производная положительна, при $50<t\le80$ отрицательна. На концах $Q(20)=Q(80)=0$, поэтому
\[
  Q_{\max}=Q(50)=30\cdot30=900.
\]

\key{Приём:} незнакомый прикладной контекст не должен пугать. Сначала выпишите данные формулы и свяжите их в одну целевую функцию.

\section*{6. Если целевая функция --- дробь}
\sectionline
В задаче на электрическую мощность могут быть заданы
\[
  I=\frac{36}{3+R},\qquad P=I^2R,\qquad R>0,
\]
где $R$ --- сопротивление нагрузки в омах. После подстановки
\[
  P(R)=\frac{36^2R}{(3+R)^2}.
\]
По правилу производной частного и последующему упрощению
\[
  P'(R)=\frac{36^2(9-R^2)}{(3+R)^4}
       =\frac{36^2(3-R)}{(3+R)^3}.
\]
В области $R>0$ знаменатель положителен, поэтому знак производной определяется множителем $3-R$:
\[
  0<R<3:~P'(R)>0,\qquad R>3:~P'(R)<0.
\]
Следовательно,
\[
  R=3~\text{Ом},\qquad
  P_{\max}=\frac{36^2\cdot3}{6^2}=108~\text{Вт}.
\]

\key{Практическое правило:} перед расстановкой знаков доводите производную до формы, где её знак читается сразу.

\clearpage
\section*{7. Геометрическая модель: коробка из квадратного листа}
\sectionline
Из квадратного листа $30\times30$ см по углам вырезают квадраты со стороной $x$ и загибают края. После вырезания сторона дна становится $30-2x$, а высота коробки равна $x$.

\begin{center}
\begin{tikzpicture}[scale=0.16,font=\small]
\def\c{4}
\tkzDefPoint(0,0){A}
\tkzDefPoint(30,0){B}
\tkzDefPoint(30,30){C}
\tkzDefPoint(0,30){D}
\tkzDrawPolygon[color=accent,line width=1.1pt](A,B,C,D)

\tkzDefPoint(\c,0){E}
\tkzDefPoint(\c,30){F}
\tkzDefPoint(30-\c,0){G}
\tkzDefPoint(30-\c,30){H}
\tkzDefPoint(0,\c){I}
\tkzDefPoint(30,\c){J}
\tkzDefPoint(0,30-\c){K}
\tkzDefPoint(30,30-\c){L}
\tkzDrawSegments[dashed,color=accent,line width=0.8pt](E,F G,H I,J K,L)

\fill[soft] (0,0) rectangle (\c,\c);
\fill[soft] (30-\c,0) rectangle (30,\c);
\fill[soft] (0,30-\c) rectangle (\c,30);
\fill[soft] (30-\c,30-\c) rectangle (30,30);

\draw[{Stealth[length=2mm]}-{Stealth[length=2mm]},darkaccent,line width=0.8pt] (0,-2.2)--(\c,-2.2) node[midway,below] {$x$};
\draw[{Stealth[length=2mm]}-{Stealth[length=2mm]},darkaccent,line width=0.8pt] (\c,15)--(30-\c,15) node[midway,fill=white] {$30-2x$};
\node at (15,-4.4) {$30$ см};
\node[rotate=90] at (-2.2,15) {$30$ см};
\end{tikzpicture}
\end{center}

Размеры коробки:
\[
  (30-2x)\times(30-2x)\times x.
\]
Поэтому
\[
  V(x)=x(30-2x)^2,
  \qquad 0<x<15.
\]
После раскрытия скобок
\[
  V(x)=900x-120x^2+4x^3,
\]
\[
  V'(x)=900-240x+12x^2=12(x-5)(x-15).
\]
Точка $x=15$ лежит на границе вырожденной конструкции и не входит в допустимую область. На $(0,5)$ имеем $V'(x)>0$, на $(5,15)$ --- $V'(x)<0$. Следовательно,
\[
  x=5~\text{см},\qquad V_{\max}=5\cdot20^2=2000~\text{см}^3.
\]

\section*{8. Частые ошибки}
\sectionline
\begin{itemize}
  \item Дифференцировать не ту величину. Сначала определите, что именно оптимизируется.
  \item Оставить в целевой формуле две или три переменные вместо одной.
  \item Найти $f'(x)=0$ и сразу объявить ответ, не доказав максимум или минимум.
  \item Исследовать производную вне допустимого промежутка.
  \item Забыть проверить границы, если они входят в область определения.
  \item Не упростить громоздкую производную перед анализом знаков.
  \item После нахождения оптимального параметра забыть вычислить само экстремальное значение, если оно требуется.
  \item Для целого количества товара округлить нецелую критическую точку без сравнения соседних допустимых целых значений.
\end{itemize}

\section*{9. Мини-памятка перед ответом}
\sectionline
Проверьте пять пунктов:
\begin{enumerate}
  \item целевая функция соответствует величине, которую нужно максимизировать или минимизировать;
  \item функция выражена через одну переменную;
  \item ограничения выписаны до анализа производной;
  \item критические и, при необходимости, граничные точки исследованы только в допустимой области;
  \item вывод о максимуме или минимуме обоснован, а ответ дан именно на вопрос задачи.
\end{enumerate}

\clearpage
\section*{Тренировочный блок --- Путь А}
\sectionline
\textbf{10 задач.} Для каждой задачи придерживайтесь полного маршрута: целевая функция $\to$ ограничения $\to$ производная $\to$ критические точки $\to$ анализ знаков и границ $\to$ вывод.

\begin{enumerate}
  \item Цена товара зависит от числа проданных единиц: $p(q)=240-4q$, где $q\in\mathbb N$. Выручка $R=pq$. Сколько единиц товара нужно продать, чтобы выручка была максимальной? Найдите максимальную выручку.

  \item Прямоугольник имеет стороны $x$ и $y$, причём $y=120-2x$, $x>0$, $y>0$. При каком $x$ площадь прямоугольника максимальна? Найдите вторую сторону и максимальную площадь.

  \item При $10\le t\le70$ заданы $A=t-10$ и $B=70-t$. Показатель $Q=AB$. Найдите значение $t$, при котором $Q$ максимален, и само максимальное значение.

  \item Сила тока задаётся формулой $I=\dfrac{48}{4+R}$, где $R>0$ измеряется в омах. Мощность на нагрузке равна $P=I^2R$. Найдите сопротивление $R$, при котором мощность максимальна, и максимальную мощность.

  \item Из квадратного листа $24\times24$ см вырезают по углам квадраты со стороной $x$ и делают открытую коробку. При каком $x$ объём коробки максимален? Найдите этот объём.

  \item Цена единицы товара равна $p(q)=200-2q$, а затраты на производство $q$ единиц составляют $C(q)=20q+500$, где $q\in\mathbb N$. При каком $q$ прибыль максимальна? Найдите максимальную прибыль.

  \item Для $x>0$ рассмотрите величину
  \[
    F(x)=x+\frac{36}{x}.
  \]
  Найдите её наименьшее значение и значение $x$, при котором оно достигается.

  \item Прямоугольник имеет площадь $144$. Его стороны положительны. При каких длинах сторон периметр прямоугольника минимален? Найдите минимальный периметр.

  \item Из прямоугольного листа $40\times30$ см по углам вырезают одинаковые квадраты со стороной $x$ и делают открытую коробку. Найдите $x$, при котором объём максимален. Ответ дайте в точном виде.

  \item Сила тока задаётся формулой $I=\dfrac{60}{5+R}$, где $R>0$ измеряется в омах, а мощность равна $P=I^2R$. Найдите $R$, при котором мощность максимальна, и максимальную мощность.
\end{enumerate}

\clearpage
\section*{Ответы}
\sectionline
\renewcommand{\arraystretch}{1.35}
\begin{longtable}{@{}>{\bfseries}p{10mm} p{0.82\linewidth}@{}}
\toprule
№ & Ответ \\
\midrule
\endfirsthead
\toprule
№ & Ответ \\
\midrule
\endhead
1 & $q=30$, $R_{\max}=3600$ ден. ед. \\
2 & $x=30$, $y=60$, $S_{\max}=1800$. \\
3 & $t=40$, $Q_{\max}=900$. \\
4 & $R=4$ Ом, $P_{\max}=144$ Вт. \\
5 & $x=4$ см, $V_{\max}=1024$ см$^3$. \\
6 & $q=45$, максимальная прибыль $3550$ ден. ед. \\
7 & $x=6$, $F_{\min}=12$. \\
8 & стороны $12$ и $12$, $P_{\min}=48$. \\
9 & $\displaystyle x=\frac{35-5\sqrt{13}}{3}$ см. \\
10 & $R=5$ Ом, $P_{\max}=180$ Вт. \\
\bottomrule
\end{longtable}

\vfill
{\small\color{muted}Лёвин Артём Александрович эксклюзивно для Николь Саркисьянц}\par

\end{document}
